\chapter{Troubleshooting}
\label{app:trouble}

\section{Building}

\begin{description}
\item[\ctp{fixed object 0 landed on pointer 0, expected 2}] You asked
  \ct{OM=bb} to bootstrap an image. It cannot---it exists to load the 1983
  Xerox image. Build with the default \ct{OM=mt}.
\item[\ctp{SDL3 not found.}] No SDL3 development package, and the build
  stops rather than quietly handing you a system with no window. Install
  \ct{SDL3-devel} or \ct{libsdl3-dev} and build again, or ask for
  \ct|make HEADLESS=1| and get everything except the window. \ct{make deps}
  reports on SDL3 and every other requirement at once, and runs on a machine
  too bare to build anything.
\item[Garbled glyphs on screen] A stale object file in the archive.
  \ct|make clean && make|.
\item[The VM says the font and the image disagree] You ran \ct{make font} and
  did not rebuild the image. Some classes fix their line grid at
  class-initialisation time, so an image built against one face cannot be shown
  another. Rebuild the image.
\end{description}

\section{The window}

\begin{description}
\item[Nothing happens when I choose \ctp{browser}] It is waiting for you to
  frame the window. Press the left button, drag out a rectangle, release. See
  Chapter~\ref{ch:tenminutes}.
\item[The menu flashes and disappears] You clicked. Press and \emph{hold}.
\item[No mouse button does anything] The display was resized without the image
  being told, so no controller believes the cursor is inside it. Run with
  \ct|ST_DISPLAY_TRACE=1| to see which step refused. As a workaround,
  \ct|ST_DISPLAY_FIT=off| keeps the image's own screen size.
\item[The screen does not update when I switch back to it] The
  window-exposed event is not arriving. \ct{restore display} on the desktop
  menu redraws everything.
\item[A window is drawn but ignores every click and keystroke] The image
  rebuilds its idea of which buttons are down from the event stream and from
  nothing else, so a lost button-release leaves it believing a button is
  still held: \ct{Sensor noButtonPressed} never answers true again, no
  controller takes control, and the window sits there looking perfectly
  fine. The gesture that risks it is dragging at full rate with a button
  held while the image is busy---framing a window, or any rubber band.
  Positions coalesce in the event ring and transitions never do, so the
  events that cannot be reconstructed are the ones that are never dropped.
  If you ever see
\begin{plain}
st80: N input events were dropped
\end{plain}
  at exit, the image's button state may have been left wrong, and that
  message is worth reporting.
\item[The rubber band blinks while I drag] The 1983 idiom draws the
  rectangle as a flash between two cancelling blits, and the band exists
  only for the gap between them---too short for a frame to sample. The drawn
  half is forced to the screen before the undo lands rather than left for
  the next frame to catch.
\item[The window is tiny / has huge borders] Set the window size with
  \ct|ST_DISPLAY_WINDOW=WxH|, and the pixel scale with
  \ct|ST_DISPLAY_SCALE|.
\end{description}

\section{Writing code}

\begin{description}
\item[\ctp{nil is not a Boolean, so it cannot decide a conditional}] A
  \ct{NonBooleanReceiver}: the receiver of an inlined \ct{ifTrue:},
  \ct{and:} or \ct{whileTrue:} was not \ct{true} or \ct{false}---almost always
  an uninitialised instance variable holding \ct{nil}. The message names the
  object that was there instead.
\item[\ctp{doesNotUnderstand: \#ifTrue:}] The same thing where the send was
  not inlined.
\item[A global is \ctp{nil} that should not be] Look in \ct{Undeclared}. A
  reference to a name nothing defines compiles quietly, holding \ct{nil}. A
  typo in a class name does exactly this.
\item[My \ctp{Set} lost an object I put in it] Either you overrode \ct{=}
  without \ct{hash}, or you mutated an object while it was an element.
\item[\ctp{3 + 4 * 2} answered 14] Binary selectors have no precedence among
  themselves. Parenthesise.
\item[A cascade went to the wrong receiver] A cascade sends to the receiver of
  the \emph{last message before the first semicolon}. Add \ct{yourself} if you
  want the receiver back.
\item[\ctp{10 / 4} answered \ctp{(5/2)}] \ct{/} answers an exact
  \ct{Fraction}. Use \ct{//} for integer division or \ct{asFloat}.
\item[\ctp{more than 63 literals referenced}] A method may reference at most
  63 literals; the header holds the count in six bits. Every distinct string,
  symbol, large integer and literal array counts, and repeated strings are not
  shared---two occurrences of \ct|'x'| are two mutable objects. Split the
  method, or move the data into a collection built once.
\item[\ctp{a method overflowed its frame}] The context was sized from the
  compiler's estimate of the working stack and the method pushed past it.
  Rare, and worth reporting along with the method that did it.
\end{description}

\section{JSON}

\begin{description}
\item[A number came back as a \ctp{Fraction}] Expected, and the point.
  \ct{'1.5' asJson} is \ct{3/2}, so that a tenth times ten is one rather than
  1.0000000000000002. Ask \ct{floatAt:} when a \ct{Float} is what you want.
  Chapter~\ref{ch:json}.
\item[My document came back wrapped in quotes] You sent \ct{asJsonString} to
  a String. That is the String \emph{written as} a JSON value, which is what
  it looks like in quotes. To read text as a document, send \ct{asJson}.
\item[The pretty-printed text is one long line on screen] It breaks lines with
  \ct{Character nl}, 10, which is what everything outside this image reads.
  A \ct{Paragraph} breaks on the carriage return, 13, because the Alto did.
  Translate where the text meets the screen; a file written through
  \ct{FileDirectory>>textFile:} is translated for you.
\item[\ctp{a value cannot begin with \$\} at character 12}] The parser is
  RFC 8259 and refuses what most libraries forgive---a trailing comma, an
  unquoted name, single quotes, a leading zero, a bare word, a second document
  after the first. The position in the message is where to look.
\item[An object with thousands of names is slow] You are running the
  \ct{bluebook} profile. 1983's \ct{String>>hash} reads the first character,
  the second-to-last, and the length; every other profile loads
  \ct{lib/Collections-Protocol}, whose \ct{String>>hash} reads them all. See
  \ct{doc/JSON.md}.
\end{description}

\section{Concurrency}

\begin{description}
\item[It works at one worker and fails at eight] It is a race. Read
  Chapter~\ref{ch:contract}, then look for an unsynchronised
  read-modify-write, a shared collection, a lazily-initialised class
  variable---or a class variable used as scratch space, which is what 1983's
  \ct{SmallInteger>>printOn:base:} did with \ct{Digitbuffer} until
  \ct{lib/Concurrency} replaced it. Nothing about the symptom said
  \emph{printing}: Strings came back with another worker's digits in them.
\item[A process is waiting forever] If it is on a \ct{Semaphore} you are using
  for mutual exclusion, you may have taken it twice. Use \ct{Mutex}, which
  detects that and raises rather than deadlocking.
\item[\ctp{every process is blocked; nothing can run}] The scheduler found
  nothing ready, no timer armed and no signal queued, and prints what each
  worker was holding. On one worker it usually means what it says. On many it
  once meant a process the scheduler had lost, and the line under it is the
  evidence to bring.
\item[The Transcript output is interleaved] Between sends, expected: lines
  from different workers arrive in the order they were sent. Within one
  send it cannot be---each \ct{show:} holds the Transcript's lock---so build
  the line in a \ct{WriteStream} and \ct{show:} it once.
\item[\ctp{Processor workerCount} answers 1] Expected: neither
  \ct{-bootstrap -eval} nor \ct{-run} starts the worker pool today. The
  parallel test suites and \ct{make bench} do.
\end{description}

\section{Saving}

\begin{description}
\item[The Browser shows no source for my methods] The image has been moved
  away from its \ct{.changes} file. Source is stored as a pointer into that
  file, not in the image.
\item[The display came back 992 wide instead of 1000] The image's own
  arithmetic: \ct{DisplayScreen} rounds a Form's width down to a sixteen-bit
  word boundary when saving shrinks and restores the display. Harmless.
\end{description}

\section{Getting more information}

\begin{center}
\small
\begin{tabular}{@{}ll@{}}
\toprule
\ctp{ST\_DISPLAY\_TRACE=1} & why a resize was refused; draw and present counts\\
\ctp{ST\_EVAL\_TRACE=1} & every bytecode an \ctp{-eval} runs\\
\ctp{ST\_BOTTOM\_LOG=1} & a return off the bottom of a stack\\
\ctp{./st80 -disasm} & what the compiler actually emitted\\
\ctp{make ASAN=1 test} & memory errors\\
\ctp{make TSAN=1 test} & data races\\
\bottomrule
\end{tabular}
\end{center}

\section{The demo application}

\begin{description}
\item[The server starts, and every screen says there is no driver] The
  SQLite ODBC driver is not installed: \ct{odbcinst -q -d} must list
  \ct{[SQLITE3]}. It is \ct{sqliteodbc} on Fedora and \ct{libsqliteodbc} on
  Debian. The demo reaches SQLite through ODBC like every database here.
\item[An https URL answers \emph{this build has no TLS}] OpenSSL was not
  found when \ct{st80} was built, or \ct{NOTLS=1} was asked for. \ct{make deps}
  says which; install \ct{openssl-devel} (Fedora) or \ct{libssl-dev}
  (Debian) and build again. \ct{Socket isTlsAvailable} answers for an image.
\item[\emph{certificate verify failed}] The server's certificate is not one
  the system's store vouches for, or is not for the name you asked for; the
  words after the colon say which---\emph{hostname mismatch},
  \emph{certificate has expired}, \emph{self-signed certificate}. There is no
  way to say ``trust it anyway'', on purpose. A private authority goes in
  the store, or in \ct{SSL_CERT_FILE} in the environment.
\item[An LLMError says \emph{answered 401}] The key was refused. It is read
  from the environment \ct{st80} was started in---\ct{Smalltalk environmentAt:}
  shows what the image sees---and given with \ct{apiKey:}.
\item[A model \emph{asked for tools twenty times in a row}] Every answer was
  another tool call: the tool's result is not what the model needs, or the
  prompt asks for something no tool answers. \ct{messages} on the
  conversation shows the exchange.
\item[The Ollama screen says the server is not running, and it is] An
  older image took the first address \ct{localhost} resolves to, which is
  often \ct{::1} while Ollama listens on \ct{127.0.0.1} alone;
  \ct{connectTo:port:} tries every address now. Rebuild the image, or name
  the server outright with \ct{OllamaUrl} in \ct{server.json}.
\item[The login page appears, then \emph{Error communicating with the server}]
  The pages are being served but nothing answers \ct{/rest}: you started
  an \ct{HttpServer} with the demo's front end as its document root, which
  serves files and no more. The demo needs a \ct{RestServer}---from a
  workspace \ct{Smalltalk at: #Demo put: (RestServer fromConfigFile: 'demo/server.json')}
  and then \ct{Demo start}, or
  \ct{./st80 -serve demo.im demo/server.json}. The same
  message means the same thing under \ct{-serve}: the server is not
  running, or is on another port than the page expects.
\item[Report and Export answer \emph{No such method}] They are not
  written: Kiss makes a PDF through \ct{groff} and serves the file back,
  and this server does not yet serve files it made.
\item[A service will not load: \ctp{... is undeclared}] A name in the
  file that is neither a temporary, an instance variable, a class variable
  nor a global. A misspelt class name is the usual one. The message names
  the identifier; compare it with what \ct{Smalltalk} holds.
\end{description}
